% Slide 9: Reference Configuration Pullback (The Diffeomorphism)
\begin{frame}{Reference Configuration Pullback (The Diffeomorphism)}
    \begin{columns}
        \begin{column}{0.6\textwidth}
            \textbf{Pullback to Static Reference Domain $\hat{\Omega}$:}
            \begin{itemize}
                \item Morphing mapping: $\vect{x} = \boldsymbol{\Psi}(\boldsymbol{\xi}; \boldsymbol{\mu})$.
                \item Integrations are performed over the static grid of $\hat{\Omega}$ where basis coordinates are fixed.
                \item Shape deformations are shifted entirely into the parametric Jacobian $\mat{J}(\boldsymbol{\xi}; \boldsymbol{\mu})$ evaluated locally at Gauss points.
            \end{itemize}
            \vspace{0.2cm}
            \textbf{Diffeomorphic Properties:}
            \begin{itemize}
                \item Smoothness ($C^{\infty}$) ensures no fold-overs or singular nodes.
                \item Preserves the vector space structure for offline POD snapshots.
            \end{itemize}
        \end{column}
        \begin{column}{0.38\textwidth}
            \centering
            \begin{tikzpicture}[scale=0.8]
                % Reference Configuration
                \draw[draw=navyblue, fill=navyblue!5, thick] (-1.5,1.2) rectangle (1.5,2.0);
                \node at (0,1.6) {\scriptsize Reference $\hat{\Omega}$ (Static)};
                
                % Pullback arrow
                \draw[->, thick, cyansec, line width=1.5pt] (0,1.1) -- (0,0.3) node[midway, right] {$\boldsymbol{\Psi}(\boldsymbol{\xi}; \boldsymbol{\mu})$};
                
                % Deformed Physical
                \shade[left color=blue!20, right color=blue!5] 
                (-1.5, -0.8) -- (-0.6, -0.8) 
                .. controls (-0.3, -0.8) and (-0.2, -0.5) .. (0, -0.5)
                .. controls (0.2, -0.5) and (0.3, -0.8) .. (0.6, -0.8)
                -- (1.5, -0.8) -- (1.5, 0.0) -- (-1.5, 0.0) -- cycle;
                \draw[draw=navyblue, thick] 
                (-1.5, -0.8) -- (-0.6, -0.8) 
                .. controls (-0.3, -0.8) and (-0.2, -0.5) .. (0, -0.5)
                .. controls (0.2, -0.5) and (0.3, -0.8) .. (0.6, -0.8)
                -- (1.5, -0.8) -- (1.5, 0.0) -- (-1.5, 0.0) -- cycle;
                \node at (0,-0.3) {\scriptsize Physical $\Omega(\boldsymbol{\mu})$};
            \end{tikzpicture}
        \end{column}
    \end{columns}
\end{frame}

% Slide 10: Mathematical Pullback of Weak Form
\begin{frame}{Mathematical Pullback of Weak Form}
    \textbf{Stiffness Assembly in Reference Space $\hat{\Omega}$:}
    \begin{equation*}
        \mat{K}_{ab}(\boldsymbol{\mu}) = \int_{\hat{\Omega}} \mat{B}_a^T(\boldsymbol{\xi}; \boldsymbol{\mu}) \mat{D} \mat{B}_b(\boldsymbol{\xi}; \boldsymbol{\mu}) \det\mat{J}(\boldsymbol{\xi}; \boldsymbol{\mu}) \, t_h \, d\hat{\Omega}
    \end{equation*}
    where the physical derivatives in B-matrix $\mat{B}_a$ are calculated via inverse Jacobian transformation:
    \begin{equation*}
        \begin{bmatrix} \frac{\partial R_a}{\partial x} \\[0.3em] \frac{\partial R_a}{\partial y} \end{bmatrix} = \mat{J}^{-1}(\boldsymbol{\xi}; \boldsymbol{\mu}) \begin{bmatrix} \frac{\partial R_a}{\partial \xi} \\[0.3em] \frac{\partial R_a}{\partial \eta} \end{bmatrix}
    \end{equation*}
    \begin{itemize}
        \item \textbf{Benefit}: The integration grid is static. The parameter-dependence is fully contained inside $\mat{J}^{-1}$ and $\det\mat{J}$.
        \item \textbf{Fixed Vector Space}: Enables SVD to extract a coordinate-consistent basis.
    \end{itemize}
\end{frame}

% Slide 11: Reduced Order Modeling: POD Subspace Construction
\begin{frame}{Reduced Order Modeling: POD Subspace Construction}
    \begin{itemize}
        \item \textbf{Proper Orthogonal Decomposition (POD)}: Find a low-dimensional basis $\boldsymbol{\Phi} \in \mathbb{R}^{N \times r}$ that best captures the snapshot space.
        \item \textbf{Snapshot Matrix}: Collect $M_{\text{snap}}$ displacement vectors under different parameter sets $\boldsymbol{\mu}^{(k)}$:
        \begin{equation*}
            \mat{S}_u = \left[ \underline{\mathbf{U}}(t_1; \boldsymbol{\mu}^{(1)}), \, \underline{\mathbf{U}}(t_2; \boldsymbol{\mu}^{(1)}), \, \dots, \, \underline{\mathbf{U}}(t_T; \boldsymbol{\mu}^{(P)}) \right]
        \end{equation*}
        \item \textbf{Singular Value Decomposition (SVD)}:
        \begin{equation*}
            \mat{S}_u = \mat{V} \boldsymbol{\Sigma} \mat{W}^T
        \end{equation*}
        \item \textbf{Subspace Dimension}: The first $r$ columns of $\mat{V}$ corresponding to the largest singular values are selected as the reduced basis $\boldsymbol{\Phi}$.
    \end{itemize}
\end{frame}

% Slide 12: Galerkin ROM & The Reconstruction Error
\begin{frame}{Galerkin ROM \& The Reconstruction Error}
    \begin{columns}
        \begin{column}{0.48\textwidth}
            \textbf{Galerkin Projection of Equation of Motion:}
            \begin{equation*}
                \boldsymbol{\Phi}^T \mat{M} \boldsymbol{\Phi} \ddot{\underline{\mathbf{q}}}(t) + \boldsymbol{\Phi}^T \mat{C} \boldsymbol{\Phi} \dot{\underline{\mathbf{q}}}(t) + \boldsymbol{\Phi}^T \mathbf{F}_{\text{int}}(\boldsymbol{\Phi}\underline{\mathbf{q}}; \boldsymbol{\mu}) = \boldsymbol{\Phi}^T \mathbf{F}_{\text{ext}}
            \end{equation*}
            \textbf{Subspace Dimension Analysis:}
            \begin{itemize}
                \item At $r=8$ modes, the relative reconstruction error falls below $0.2\%$.
                \item At $r=15$ modes, error converges to $0.012\%$.
            \end{itemize}
            \vspace{0.2cm}
            \textbf{The Remaining Bottleneck:}
            \begin{itemize}
                \item Evaluating $\boldsymbol{\Phi}^T \mathbf{F}_{\text{int}}$ still requires integrating over all $3,200$ elements online ($\mathcal{O}(N)$ scaling).
            \end{itemize}
        \end{column}
        \begin{column}{0.48\textwidth}
            \centering
            \begin{tikzpicture}[scale=0.8]
                \begin{axis}[
                    width=1.0\textwidth,
                    height=0.75\textwidth,
                    xlabel={Number of POD Modes $r$},
                    ylabel={Relative $L_2$ Error (\%)},
                    xmin=1, xmax=10,
                    ymin=0.0, ymax=2.5,
                    grid=both,
                    grid style={line width=.1pt, draw=gray!10},
                    major grid style={line width=.2pt, draw=gray!20},
                    axis lines=left,
                    xtick={1,2,3,4,5,6,7,8,9,10},
                    ytick={0.0,0.5,1.0,1.5,2.0,2.5},
                    x tick label style={font=\tiny},
                    y tick label style={font=\tiny},
                ]
                    \addplot[color=cyansec, thick, mark=*] coordinates {
                        (1, 2.3424) (2, 1.7880) (3, 1.3420) (4, 0.9011) 
                        (5, 0.6542) (6, 0.4010) (7, 0.2874) (8, 0.1985) 
                        (9, 0.1432) (10, 0.1062)
                    };
                \end{axis}
            \end{tikzpicture}
        \end{column}
    \end{columns}
\end{frame}
